% Chapter 4: Simulation and Results
% Every number below is listed with its source in thesis/numbers.md and is regenerated by
% thesis/make_results.py.
\providecommand{\tbd}[1]{\textbf{[TBD: #1]}}

\chapter{Simulation and Results}
\label{ch:results}

This chapter evaluates the framework of Chapter~3 in two stages. In Stage~1 we replicate the
method of Haugh and Singal~\cite{haugh2021} on the FantasyGo round contests of the 2023-24
season, which are the pool-style setting for which the method was designed; this stage
validates our implementation. In Stage~2 we apply the method to the 2024-25 Fantasy Premier
League (FPL) head-to-head (H2H) league, which is the contribution of this study. In both stages
the mean-variance (MVT) optimizer is compared with a baseline optimizer that maximises
predicted points and ignores the opponents, and in both stages we study how the benchmark
rank $r$ changes the chosen teams and their outcomes.

\section{Experimental Design}
\label{sec:design}

\subsection{Data and information set}
\label{sec:design-data}

\paragraph{Stage 1 (2023-24).} The player data are the FPL player-gameweek records of the
2023-24 season. The contest data are the FantasyGo round contests of gameweeks 6 to 38: for
each of the 33 gameweeks, every entry's team, captain, vice-captain, points and prize. A contest
had between 203 and 342 entries. Predictions start at gameweek 5, so gameweek 6 is the first
with one gameweek of prediction errors.

\paragraph{Stage 2 (2024-25).} The player data are the FPL player-gameweek records of the
2024-25 season, and the league data are the weekly 15-man squads of the 20 managers of one FPL
H2H league. Our own team is manager 205, and the fixtures are a single round-robin in each half
of the season. The rolling season runs from gameweek 6 to 38, again 33 gameweeks.

\paragraph{Information set.} Every decision for gameweek $t$ uses only data available before
its deadline: features, prices and ownership up to gameweek $t-1$, the selection models fitted
on gameweeks before $t$, and prediction errors of gameweeks before $t$. Teams are scored on the
actual points of gameweek $t$.

\paragraph{Gameweeks without a target.} The source data have no expected points (xP) in
gameweek 26 of 2023-24 and in gameweeks 22, 32 and 34 of 2024-25: xP is recorded as zero there,
while the players' points are real. These gameweeks are left out of every xP metric (model
training, mean absolute error, predicted-versus-actual comparisons) and kept in every points and
win metric.

\subsection{H2H as a double-up, and the method of Haugh and Singal}
\label{sec:design-method}

We use the formulation and solution procedure of~\cite{haugh2021} without change to their
logic. Our team $w \in \mathbb{W}$ is chosen to solve problem (3.16),
\[
  \max_{w\in\mathbb{W}} \; \mathrm{Prob}\bigl\{ w^{T}\delta > G^{(r')}(\mathbb{W}_o,\delta) \bigr\},
\]
where $\delta \sim N(\mu_\delta,\Sigma_\delta)$ are the players' points and $G^{(r')}$ is the
$r'$-th order statistic of the scores $G_o = w_o^{T}\delta$ of the teams $w_o\in\mathbb{W}_o$.
Theorem~3.2.3.1 reduces (3.16) to the mean-variance problems (3.17), one branch for
$\mu_{Y_w}\ge 0$ and one for $\mu_{Y_w}<0$; the moments $\mu_{G^{(r')}}$,
$\sigma^2_{G^{(r')}}$ and $\sigma_{\delta,G^{(r')}}$ are estimated from Monte Carlo samples
of $(\delta,G^{(r')})$ generated as in Lemma~3.2.2; and Algorithm~1 solves a mixed-integer
quadratic program (MIQP) for every $\lambda$ in a grid $\Lambda$ and returns the team $w(\lambda^*)$
with $\lambda^* = \arg\max_{\lambda\in\Lambda}\widehat{\mathrm{Prob}}\{Y_{w(\lambda)}>0\}$.
Two details of the implementation, common to both stages, make Algorithm~1 concrete. The branch
of (3.17) is set once per gameweek by the sign of $\mu_{Y_w}$ for a reference team (the
$\lambda=0$ team in Stage~1, our current squad in Stage~2). And $\Sigma_\delta$ enters the MIQP
through its eigendecomposition, so that the variance is a sum of squares, which the solver SCIP
handles far faster than the pairwise form.

In a pool contest $\mathbb{W}_o$ is the field of other entries, and $r$ is the share of the
field we must beat: $r=50\%$ is a double-up, and $r=100\%$ is winner-takes-all. In H2H there is
a single opponent, and $\mathbb{W}_o$ becomes the set of squads that this opponent could
plausibly field in the coming gameweek (§3.1.3). Beating the opponent's median plausible squad
($r=50\%$) is then the H2H analogue of the double-up, and beating their best plausible squad
($r=100\%$) is the analogue of winner-takes-all. In both stages the rank is taken from the top:
$G^{(r')}$ is the score with index $\lceil (1-r)\,O\rceil$ from the top, at least $1$, among the
$O$ teams. In Stage~1, $O$ counts every simulated entry, duplicates included, because a real
field contains duplicates too. In Stage~2, $\mathbb{W}_o$ is deduplicated and $O$ counts
distinct squads, so $r$ is a rank among the distinct squads the opponent could plausibly field
rather than a probability level. The two definitions coincide at $r=100\%$, where $G^{(r')}$ is
the maximum.

\paragraph{Non-strategic opponents.} As in~\cite{haugh2021}, the opponents are assumed not to
react to our selection: $\mathbb{W}_o$ is built only from their own historical selection
behaviour, and our team does not enter its construction. We return to this assumption in
Section~\ref{sec:discussion}.

\subsection{The baseline optimizer}
\label{sec:design-baseline}

In each stage the comparison is between MVT and a baseline that applies the optimizer directly
to the player data and the predicted xP. Each gameweek the baseline picks the team and captain
with the highest total predicted xP under that season's rules, with no opponent model, no
$\mathbb{W}_o$ and no Monte Carlo, and it is scored on the actual points. The baseline does not
depend on $r$, so one baseline season serves every $r$. The real managers and the real opponents
are shown as reference points only.

\paragraph{Points distribution.} In both stages $\mu_\delta$ is the predicted xP of every player
with a fixture. $\Sigma_\delta$ is the sample covariance of past prediction errors over pairs of
players with at least three gameweeks in common, with each player's own variance on the
diagonal. A player with too short a history receives the median variance of his position and no
covariances; before any player has three gameweeks of errors, the variance of his position is
pooled over every past player-gameweek. The correlation matrix is made positive definite by
flooring its eigenvalues at $10^{-4}$ of the largest, and rescaled to a unit diagonal. The
errors differ between the stages and are described there.

\section{Player Performance Predictions}
\label{sec:predictions}

Both optimizers take their mean points $\mu_\delta$ from a machine-learning model of each
player's xP in the next gameweek, and $\Sigma_\delta$ from that model's past errors. This
section describes the models and the walk-forward procedure that chooses them, and reports
their accuracy in both seasons. The two seasons use the same candidate models and the same
procedure. They differ in their features and in the players the models are fitted and scored
on.

\subsection{Target and features}
\label{sec:predictions-features}

The target is FPL's expected points (xP) for the gameweek, not the realised points; xP is less
noisy, and in Stage~2 the noise of the realised points around xP enters $\Sigma_\delta$
separately (Section~\ref{sec:s2-setting}). A separate model is fitted for each position. Every model uses
three categorical features (the player's club, the opponent, and whether the match is at home)
and a set of numeric features of four kinds:
\begin{itemize}
  \item \emph{fixture strength}: the Elo ratings of the club and the opponent, FPL's fixture
  difficulty ratings and the win probability;
  \item \emph{market}: the player's price, the share of FPL managers who own the player,
  and the week's transfers in and out;
  \item \emph{recent form}: averages over the previous three gameweeks of xP, minutes, starts,
  the ICT index and its influence, creativity and threat components, and the underlying
  statistics relevant to the position (expected goals, assists and goal involvements for
  outfield players; expected goals conceded, clean sheets and saves for goalkeepers and
  defenders);
  \item \emph{season-specific}: xGChain and xGBuildup in 2023-24, which the 2024-25 data source
  does not contain, and tackles in 2024-25 (except for forwards), which are not published for
  2023-24.
\end{itemize}
The 2023-24 models have 25 numeric features (28 for defenders), and the 2024-25 models between
20 and 24. All features of gameweek $t$ are known before its deadline.

\subsection{Candidate models and walk-forward selection}
\label{sec:predictions-selection}

Eight base learners compete: ridge regression ($\alpha=1$), a decision tree (depth 6, at least
30 rows per leaf), a random forest (300 trees, at least 10 rows per leaf), histogram gradient
boosting (scikit-learn defaults), LightGBM (300 trees, learning rate 0.03, 15 leaves), XGBoost
(300 trees, learning rate 0.03, depth 3), CatBoost (500 iterations, learning rate 0.03, depth 4)
and a support vector regressor with an RBF kernel ($C=1$). The ridge and SVR models receive
median-imputed, standardised numeric features and one-hot categories; the tree models receive
the raw features and ordinal-coded categories. The hyperparameters are fixed in advance and
not tuned, since a season gives too few gameweeks to tune them out of sample. Two ensembles
complete the candidates: the mean of five base learners (ridge, histogram gradient boosting,
LightGBM, XGBoost and SVR), and the mean of the three base learners with the lowest MAE on the
three gameweeks before the selection point. A constant predictor, the training mean, is kept as
a floor and never selected.

The candidates are compared at fixed selection points: gameweeks 5, 10, 15, 20, 25, 30 and 35
in 2023-24, and 5, 10, 15, 21, 26, 31 and 36 in 2024-25. At each point every candidate is
trained on all earlier gameweeks and validated on that gameweek, and the one with the lowest
validation mean absolute error (MAE) is used, unchanged, until the next point, that is, for up
to five gameweeks. A selection point's own gameweek is excluded from the accuracy figures below,
since it was used to choose the model. The gameweeks without an xP target are excluded from
training and scoring.

\paragraph{Which players.} The 2023-24 models are trained and validated only on player-gameweeks
in which the player played. The 2024-25 models are trained on the players with minutes in the
previous three gameweeks, which is known before the deadline, and validated on every player in
the pool. Both are then used to predict every player with a fixture, because the optimizers
choose from the whole pool. About 40\% of the pool has an xP of zero in any gameweek, mostly
players who do not play. The two seasons' headline accuracy figures are therefore measured on
different populations, and we report both populations for both seasons below.

\subsection{Accuracy}
\label{sec:predictions-accuracy}

Table~\ref{tab:predictor-mae} reports the accuracy of the procedure on the gameweeks it was
used for, gameweeks 6 to 38 without the gameweeks lacking a target (32 gameweeks in 2023-24, 30
in 2024-25), each on the population its models were selected on. In both seasons the procedure
is well below the constant floor for every position. It is also close to, and often better
than, the best single candidate fixed in hindsight, which is not available at decision time:
in 2023-24 it is within 0.03 of it in every position and better for goalkeepers and defenders,
and in 2024-25 it is better in every position except forwards (0.743 against 0.739). The
goalkeeper model of 2023-24 is the weakest ($R^2=0.383$).

\begin{table}[ht]
  \centering
  \caption{Walk-forward accuracy of the xP predictions on the gameweeks each selected model was
  used for (against xP), on the population each season's models were selected on: players who
  played in 2023-24, every player in 2024-25. ``Best fixed'' is the single candidate with the
  lowest MAE over the same gameweeks, known only in hindsight.}
  \label{tab:predictor-mae}
  \begin{tabular}{llrrrr}
    \toprule
    Season & Position & MAE & $R^2$ & Best fixed MAE & Constant MAE \\
    \midrule
    2023-24 & GK  & 1.045 & 0.383 & 1.055 & 1.448 \\
            & DEF & 0.912 & 0.608 & 0.915 & 1.560 \\
            & MID & 0.874 & 0.665 & 0.870 & 1.633 \\
            & FWD & 1.109 & 0.573 & 1.081 & 1.853 \\
    \midrule
    2024-25 & GK  & 0.617 & 0.693 & 0.669 & 2.160 \\
            & DEF & 0.590 & 0.722 & 0.591 & 1.622 \\
            & MID & 0.594 & 0.762 & 0.597 & 1.725 \\
            & FWD & 0.743 & 0.723 & 0.739 & 2.194 \\
    \bottomrule
  \end{tabular}
\end{table}

\paragraph{Comparing the seasons.} The lower MAE of 2024-25 in Table~\ref{tab:predictor-mae}
is largely an effect of the population. Table~\ref{tab:predictor-pop} scores the deployed
predictions of both seasons on both populations. On the players who played, 2024-25 is more
accurate by 0.08 to 0.11 points for outfield positions and 0.05 for goalkeepers, and its $R^2$
is similar to 2023-24's (higher for outfield players, lower for goalkeepers). Across the whole
pool the gap is larger, for a different reason: the 2023-24 models never saw a player who did
not play, and they over-predict such players. For goalkeepers, where most of the pool are
reserves, their $R^2$ over the whole pool is $-0.05$. The 2024-25 models, trained on recently
active players and validated on the whole pool, predict the non-players close to their xP of
zero. The optimizers use the predictions for the whole pool, so in 2023-24 they see inflated
predictions for non-players; the top decile, from which they choose, is nevertheless close to
unbiased against xP (below).

\begin{table}[ht]
  \centering
  \caption{Accuracy of the deployed predictions on every player with a fixture, and on the
  players who played. Rows per gameweek are means.}
  \label{tab:predictor-pop}
  \small
  \begin{tabular}{llrrrrrr}
    \toprule
    & & \multicolumn{3}{c}{All players} & \multicolumn{3}{c}{Players who played} \\
    \cmidrule(lr){3-5}\cmidrule(lr){6-8}
    Season & Position & Rows/GW & MAE & $R^2$ & Rows/GW & MAE & $R^2$ \\
    \midrule
    2023-24 & GK  &  88 & 1.424 & $-0.052$ &  20 & 1.053 & 0.421 \\
            & DEF & 250 & 0.901 & 0.526    &  98 & 0.913 & 0.596 \\
            & MID & 334 & 0.891 & 0.631    & 136 & 0.876 & 0.664 \\
            & FWD & 100 & 0.933 & 0.606    &  37 & 1.112 & 0.574 \\
    \midrule
    2024-25 & GK  &  75 & 0.612 & 0.694    &  20 & 0.998 & 0.352 \\
            & DEF & 242 & 0.587 & 0.722    &  99 & 0.825 & 0.619 \\
            & MID & 323 & 0.593 & 0.757    & 149 & 0.798 & 0.671 \\
            & FWD &  79 & 0.735 & 0.722    &  33 & 1.006 & 0.635 \\
    \bottomrule
  \end{tabular}
\end{table}

\paragraph{The candidates.} Table~\ref{tab:predictor-candidates} gives the MAE of every
candidate as if it had been used throughout. Gradient boosting and the ensembles lead in both
seasons, and the single decision tree, ridge regression and the constant predictor trail. Of
the 28 choices per season (seven points, four positions), 2023-24 chose a gradient-boosting
model 14 times (CatBoost 7, LightGBM 5, histogram boosting 2), an ensemble 10 times, a random
forest twice and ridge and SVR once each. 2024-25 chose an ensemble 14 times (the five-model mean
10 times), gradient boosting 12 times (LightGBM 7, histogram boosting 5) and a random forest
twice. No candidate is best everywhere: CatBoost is the best fixed choice for goalkeepers and
defenders in 2023-24 but among the weaker candidates in 2024-25. Since the ranking changes with
the position and the season, and from one selection point to the next, choosing the model by
validation is safer than committing to one in advance, and the procedure's results in
Table~\ref{tab:predictor-mae} bear this out.

\begin{table}[ht]
  \centering
  \caption{MAE of each candidate over the deployed gameweeks if it had been used throughout,
  on the population of Table~\ref{tab:predictor-mae}. ``Procedure'' is the walk-forward
  selection.}
  \label{tab:predictor-candidates}
  \small
  \setlength{\tabcolsep}{4pt}
  \begin{tabular}{lrrrrrrrr}
    \toprule
    & \multicolumn{4}{c}{2023-24} & \multicolumn{4}{c}{2024-25} \\
    \cmidrule(lr){2-5}\cmidrule(lr){6-9}
    Candidate & GK & DEF & MID & FWD & GK & DEF & MID & FWD \\
    \midrule
    Procedure          & 1.045 & 0.912 & 0.874 & 1.109 & 0.617 & 0.590 & 0.594 & 0.743 \\
    \midrule
    Mean of five       & 1.063 & 0.921 & 0.870 & 1.084 & 0.669 & 0.591 & 0.607 & 0.750 \\
    Top three          & 1.072 & 0.919 & 0.874 & 1.081 & 0.719 & 0.597 & 0.606 & 0.766 \\
    CatBoost           & 1.055 & 0.915 & 0.870 & 1.091 & 0.837 & 0.635 & 0.650 & 0.796 \\
    LightGBM           & 1.118 & 0.917 & 0.888 & 1.137 & 0.721 & 0.599 & 0.599 & 0.739 \\
    XGBoost            & 1.080 & 0.926 & 0.889 & 1.115 & 0.726 & 0.615 & 0.608 & 0.758 \\
    Hist.\ boosting    & 1.093 & 0.955 & 0.908 & 1.189 & 0.763 & 0.616 & 0.597 & 0.748 \\
    Random forest      & 1.070 & 0.954 & 0.898 & 1.129 & 0.916 & 0.649 & 0.627 & 0.775 \\
    SVR                & 1.062 & 0.994 & 0.923 & 1.121 & 1.027 & 0.654 & 0.700 & 0.850 \\
    Ridge              & 1.189 & 1.032 & 0.971 & 1.229 & 0.814 & 0.705 & 0.693 & 0.994 \\
    Decision tree      & 1.139 & 1.054 & 1.012 & 1.235 & 0.949 & 0.745 & 0.733 & 0.863 \\
    Constant           & 1.448 & 1.560 & 1.633 & 1.853 & 2.160 & 1.622 & 1.725 & 2.194 \\
    \bottomrule
  \end{tabular}
\end{table}

\paragraph{Bias for the top players.} Both optimizers choose the players with the highest
$\mu_\delta$, so the bias of the predictions for those players matters more than the average
error. Among the top decile of predictions in each gameweek, the predictions are nearly unbiased
for xP (mean prediction minus xP $-0.053$ in 2023-24 and $-0.117$ in 2024-25), but they exceed
the realised points by about one point per player ($+1.009$ and $+0.976$). The highly predicted
players score, on average, less than their xP. Any team built from the top of the predictions,
the baseline's and MVT's alike, therefore expects more points than it will score, and the
estimated probabilities $\widehat{\mathrm{Prob}}\{Y_w>0\}$ inherit that optimism. Below the top
decile the 2023-24 predictions exceed xP by 0.629 points on average, against 0.084 in 2024-25,
which is the over-prediction of non-playing players described above.

\section{Stage 1: Replicating Haugh and Singal on FantasyGo}
\label{sec:s1}

\subsection{Setting and departures from the original}
\label{sec:s1-setting}

A FantasyGo round contest is a pool contest in the sense of~\cite{haugh2021}: each gameweek,
every entry picks a fresh team, and prizes are paid by rank in the field. The team rules differ
from the 15-man squad of §3.1.2. A team is 11 players (1 goalkeeper, 3 to 5 defenders, 2 to 5
midfielders, 1 to 3 forwards) within a budget of 80 at the players' FPL prices, with at most
three players from one club, a captain whose points count twice and a vice-captain whose points
count 1.5 times, and unlimited changes every gameweek. Accordingly $w$ is a vector of point
multipliers (1, 1.5 or 2 for selected players, 0 otherwise) rather than a binary vector, and
there are no bench, transfers or free-transfer constraints.

Our implementation departs from~\cite{haugh2021} in the following points, all of which are
documented in the code.
\begin{enumerate}
  \item \textbf{The field $\mathbb{W}_o$.} We simulate as many teams as there were entries in the
        previous gameweek. Each team is drawn from a league-wide Dirichlet-multinomial selection
        model of player choice per position (§3.2.1, fitted on the FantasyGo managers' earlier
        picks), with a formation drawn from the managers' earlier formations, a budget floor
        at the 25\% quantile of the real teams' earlier costs, occasional stacking of teammates
        (probability $0.35$), and the two highest-predicted players as captain and vice-captain.
        Duplicate teams are kept. The original model of~\cite{haugh2021} uses Dirichlet
        regression for the same purpose.
  \item \textbf{$\Sigma_\delta$ from points errors.} The errors are realised points minus
        predicted xP, because FantasyGo is scored on points; this includes the risk that a player
        does not play.
  \item \textbf{Variance for short histories.} Players without enough error history receive
        their position's median variance (Section~\ref{sec:design-baseline}) rather than a
        variance near zero, which would make them look riskless.
  \item \textbf{Shortlist.} The MIQP is solved over the 30 best-predicted players per position.
  \item \textbf{Solver settings.} The grid is
        $\Lambda=\{0, 0.005, 0.01, 0.02, 0.05, 0.1, 0.2, 0.5\}$, with 2{,}000 Monte Carlo
        draws per gameweek and a time limit of 30 seconds per MIQP; on the nonconvex branch
        ($\mu_{Y_w}<0$) the solver returns the best team found within the limit.
\end{enumerate}

\paragraph{The baseline.} The baseline picks, each gameweek, the 11 players, captain and
vice-captain with the highest total predicted xP under the same rules, over every player with a
fixture. Over gameweeks 6 to 38 it scores 1{,}913.5 points. The $\lambda=0$ team inside the MVT
pipeline solves the same problem over the shortlist and scores the same season total, 1{,}913.5
points, which confirms that the shortlist did not cost the baseline any points.

\subsection{The selection model and the simulated field}
\label{sec:s1-field}

The Dirichlet-multinomial selection model was refitted for every gameweek and position, 132
fits in all, each with between 2 and 8 covariates. Every fit converged (maximum
$\hat R = 1.00$, minimum bulk effective sample size 1{,}866). Table~\ref{tab:s1-oos} compares its
out-of-sample predictions of the real field's player shares with three simple alternatives. The
model has the lowest log loss for goalkeepers, defenders and forwards. For midfielders, FPL's
global ownership is marginally better (3.628 against 3.650).

\begin{table}[ht]
  \centering
  \caption{Stage 1: out-of-sample log loss of the predicted player shares in the real FantasyGo
  field, gameweeks 6--38 (lower is better).}
  \label{tab:s1-oos}
  \begin{tabular}{lrrrr}
    \toprule
    Predictor & GK & DEF & MID & FWD \\
    \midrule
    Selection model         & 2.545 & 3.792 & 3.650 & 2.769 \\
    FPL ownership           & 3.073 & 3.963 & 3.628 & 2.920 \\
    Previous gameweek's field & 5.173 & 5.429 & 4.799 & 4.059 \\
    Uniform                 & 3.268 & 4.919 & 5.204 & 3.912 \\
    \bottomrule
  \end{tabular}
\end{table}

Table~\ref{tab:s1-field} compares the simulated field with the real contest and with the
previous gameweek's real teams, all scored on the actual points. The simulated teams cost more
than the real ones (77.5 against 75.6) and score fewer points on average (52.3 against 55.5);
the score needed for the top 20\% is lower than in the real field (63.3 against 67.0). The rank
correlation between how often each player is picked in the simulated and the real field is
0.57, against 0.31 for the previous gameweek's teams. The simulated field is therefore closer to
the real one than last week's field, but it is somewhat weaker, which makes the benchmark
$G^{(r')}$ easier to beat in simulation than in the contest.

\begin{table}[ht]
  \centering
  \caption{Stage 1: the simulated field against the real contest, means over gameweeks 6--38.
  ``Top-20\% cut'' is the score of the entry at $r=80\%$.}
  \label{tab:s1-field}
  \begin{tabular}{lrrrr}
    \toprule
    Field & Cost & Mean points & Top-20\% cut & Share correlation \\
    \midrule
    Simulated           & 77.51 & 52.28 & 63.32 & 0.57 \\
    Real contest        & 75.64 & 55.51 & 66.95 & --- \\
    Previous gameweek   & 73.37 & 49.92 & 61.13 & 0.31 \\
    \bottomrule
  \end{tabular}
\end{table}

\subsection{MVT against the baseline, and the effect of $r$}
\label{sec:s1-results}

The MVT season was run at $r\in\{50,60,80,90,95,99,100\}\%$ against the simulated field.
Table~\ref{tab:s1-r} reports, for each $r$, the choice of $\lambda^*$, the risk of the chosen
team, and its results in the real contest; Figure~\ref{fig:s1-r} shows the same quantities.
A gameweek counts as a hit when the team finished in the top $1-r$ of the real contest, so the
hit target is stricter as $r$ grows (at $r=100\%$ a hit is a contest win).

\begin{table}[ht]
  \centering
  \caption{Stage 1: MVT at each $r$ against the baseline, 33 gameweeks. $\sigma_w$ is the
  standard deviation of the chosen team's points under $\Sigma_\delta$; the baseline's is 18.64.
  Hits: gameweeks finishing in the top $1-r$ of the real contest. $\Delta$ points: mean
  difference per gameweek (MVT minus baseline) with a 95\% bootstrap interval; $p$: Wilcoxon
  signed-rank test, $n=33$.}
  \label{tab:s1-r}
  \small
  \begin{tabular}{rrrrrrrrrr}
    \toprule
    $r$ & mean $\lambda^*$ & $\lambda^*\neq0$ & $\sigma_w$ & Points & $\Delta$ points [95\% CI] & $p$ & Hits (base) & Prize (base) \\
    \midrule
    50\%  & 0.114 & 32 & 14.16 & 1945.0 & $+0.95$ [$-2.98$, $4.97$] & 0.837 & 17 (13) & 7955.61 \\
    60\%  & 0.152 & 33 & 14.35 & 1952.5 & $+1.18$ [$-3.02$, $5.65$] & 0.822 & 13 (11) & 9841.08 \\
    80\%  & 0.095 & 32 & 14.55 & 1967.5 & $+1.64$ [$-2.23$, $5.59$] & 0.562 & 7 (3)   & 173.10 \\
    90\%  & 0.020 & 24 & 17.06 & 1887.0 & $-0.80$ [$-3.14$, $1.55$] & 0.409 & 3 (3)   & 123.60 \\
    95\%  & 0.004 & 16 & 20.49 & 1834.5 & $-2.39$ [$-4.70$, $-0.41$] & 0.015 & 1 (1) & 1303.81 \\
    99\%  & 0.008 & 29 & 25.46 & 1875.5 & $-1.15$ [$-4.24$, $1.94$] & 0.545 & 1 (1)   & 4490.45 \\
    100\% & 0.017 & 32 & 27.97 & 1740.5 & $-5.24$ [$-11.62$, $0.23$] & 0.042 & 1 (0) & 4536.96 \\
    \midrule
    Baseline & --- & --- & 18.64 & 1913.5 & --- & --- & --- & 1303.81 \\
    \bottomrule
  \end{tabular}
\end{table}

\begin{figure}[ht]
  \centering
  \includegraphics[width=\textwidth]{figures/s1_r_effect.pdf}
  \caption{Stage 1: the effect of $r$. (a) Season points; (b) season prize money; (c) mean
  standard deviation of the chosen team's points; (d) share of gameweeks in the top $1-r$ of the
  real contest, with MVT's mean estimated $\mathrm{Prob}\{Y_w>0\}$.}
  \label{fig:s1-r}
\end{figure}

\paragraph{The branch of (3.17) and the variance of the team.} The effect of $r$ runs through
the branch of Theorem~3.2.3.1. For $r\le 80\%$ the $\lambda=0$ team's mean exceeded
$\mu_{G^{(r')}}$ in every gameweek, so every week was on the convex branch, and a positive
$\lambda$ penalises variance: the chosen team's standard deviation fell from 18.64 points for
the baseline to between 14.16 and 14.55. At $r=90\%$ one gameweek was on the nonconvex branch,
at $r=95\%$ fifteen, and at $r=99\%$ and $100\%$ all 33. There a positive $\lambda$ rewards
variance, and the chosen team's standard deviation rose to 20.49, 25.46 and 27.97 points.
MVT thus does what Theorem~3.2.3.1 prescribes: it buys safety when it expects to be ahead of the
benchmark and buys variance when it expects to be behind. The size of $\lambda^*$ is not
comparable across the two branches, because the variance enters with opposite signs; on the
nonconvex branch the smallest positive values of $\Lambda$ already change the team.

The nonconvex MIQP did not always finish within the 30-second limit. The share of solves
stopped at the limit rises with $r$, from 56 of 264 at $r=50\%$ to 223 of 264 at $r=99\%$ and
$100\%$. At low $r$ only the values $\lambda\ge 0.1$ hit the limit, but these include the
values most often chosen as $\lambda^*$ at $r=50\%$ ($0.1$ and $0.2$), so some chosen teams are
the best found rather than proven optimal on either branch.
At high $r$ the frontier is therefore built from good, but not provably optimal, teams.

\paragraph{Points.} For $r\le 80\%$ MVT scored slightly more points than the baseline (1945.0
to 1967.5 against 1913.5), but the paired differences are small relative to their spread, and
none is significant (Wilcoxon $p\ge 0.56$). For $r\ge 90\%$ MVT scored fewer points. The
difference is significant at the 5\% level at $r=95\%$ ($-2.39$ points per gameweek, $p=0.015$)
and at $r=100\%$ ($-5.24$, $p=0.042$), but neither survives a Bonferroni correction for the
seven values of $r$ (threshold $0.007$). The loss of mean points at high $r$ is the price of the
variance that MVT deliberately buys on the nonconvex branch.

\paragraph{Hits and prizes.} At $r=50\%$ MVT finished in the top half of the contest in 17
gameweeks against the baseline's 13, and at $r=80\%$ in the top fifth in 7 against 3. Neither
difference is significant (sign test on the discordant gameweeks, $p=0.42$ and $p=0.22$). At
$r\ge 90\%$ both optimizers hit the target at most three times in 33 gameweeks, and MVT won one
contest at $r=100\%$, which the baseline never did. Prize money is dominated by single
gameweeks: of MVT's 7{,}955.61 at $r=50\%$, 7{,}832.38 came from gameweek 24, and of 9{,}841.08 at
$r=60\%$, 9{,}692.98 from gameweek 23; the baseline earned 1{,}204.85 of its 1{,}303.81 in
gameweek 26. With a handful of paid gameweeks per season, prize totals cannot separate the two
optimizers, and we do not test them.

\paragraph{Calibration.} MVT's estimates of $\mathrm{Prob}\{Y_w>0\}$ are far too high.
Table~\ref{tab:s1-calibration} pools the 231 team-gameweeks over every $r$ and compares the
estimate with the observed hit rate in the real contest. In the 69 cases where MVT estimated a
probability of at least 0.95 (mean 0.994), the team hit its target in 43\% of them; where it
estimated between 0.2 and 0.4 it hit in 2\%. At $r=50\%$ the mean estimate was 0.996, against an
observed 17 of 33 (Brier score 0.483). Three causes contribute, each identified above: the
simulated field is weaker than the real one (Table~\ref{tab:s1-field}), $\mu_\delta$ overstates
the points of the top players that MVT selects, and $\Sigma_\delta$ rests on few gameweeks of
errors. The ordering of the estimates is nevertheless informative: the observed rate rises
with the estimate above 0.4.

\begin{table}[ht]
  \centering
  \caption{Stage 1: calibration of MVT's estimated $\mathrm{Prob}\{Y_w>0\}$, all $r$ pooled
  (231 team-gameweeks).}
  \label{tab:s1-calibration}
  \begin{tabular}{lrrr}
    \toprule
    Estimate & $n$ & Mean estimate & Observed hit rate \\
    \midrule
    $[0.00, 0.20)$ & 3  & 0.189 & 0.333 \\
    $[0.20, 0.40)$ & 53 & 0.300 & 0.019 \\
    $[0.40, 0.60)$ & 39 & 0.499 & 0.000 \\
    $[0.60, 0.80)$ & 34 & 0.680 & 0.118 \\
    $[0.80, 0.95)$ & 33 & 0.865 & 0.212 \\
    $[0.95, 1.00)$ & 69 & 0.994 & 0.435 \\
    \bottomrule
  \end{tabular}
\end{table}

\paragraph{Summary of Stage 1.} The implementation reproduces the qualitative behaviour of the
method of~\cite{haugh2021}: the branch of (3.17) follows $r$, MVT reduces variance at low $r$
and increases it at high $r$, and the deliberate variance at high $r$ costs mean points. In a
single season of 33 contests none of the gains in points or hits at low $r$ is significant,
and the absolute probability estimates are not calibrated against the real contest.

\section{Stage 2: The H2H League}
\label{sec:s2}

\subsection{Setting}
\label{sec:s2-setting}

Stage~2 uses the full FPL rules of §3.1.2: a 15-man squad (2 goalkeepers, 5 defenders, 5
midfielders, 3 forwards) within the squad's value plus the money in the bank, at most three
players per club, a starting XI in a valid formation with a doubled captain, and free transfers
that are banked from week to week. The rolling season starts from our real squad before
gameweek 6 and then carries the simulated squad, free-transfer bank and money in the bank
forward. Each gameweek is scored as a real H2H match: the starting XI with the captain doubled,
with FPL's automatic substitutions and vice-captain rule, against the real opponent's real
points. The data contain squads but not line-ups, so the XI and captain of every squad not
chosen by the MIQP (the opponent's sampled squads, a held squad, the real squads) are picked by
the same rule: the highest predicted xP in a valid formation.

\paragraph{Points distribution.} The predictor targets xP, so its residuals are xP minus
prediction. These residuals form the covariance part of $\Sigma_\delta$, built as in
Section~\ref{sec:design-baseline}. A gameweek is won on points, however, and points scatter
around xP by about three times as much as xP around the prediction. The variance of realised
points about xP, pooled per position over past player-gameweeks, is therefore added to the
diagonal of $\Sigma_\delta$. The gameweeks without an xP target drop out of the residuals.

\paragraph{The two optimizers.} The baseline solves the $\lambda=0$ problem each week, which is
the convex branch whatever the opponent: the squad, XI and captain with the highest predicted
xP within the budget and free transfers. It transfers when the new squad adds more than 0.01
predicted points to holding. MVT solves the MIQP over the grid $\Lambda=\{0\}\cup\{$ten values
spaced geometrically from 0.01 to 5$\}$, with 20{,}000 opponent squad draws and 2{,}000 Monte
Carlo draws of $(\delta,G^{(r')})$, and transfers only if the best squad's estimated
$\mathrm{Prob}\{Y_w>0\}$ exceeds that of holding by 0.03; otherwise it holds, and $\lambda^*$
is undefined for that week. Both use the same seed (42). For all seven MVT seasons the MIQP time
limit is 60 seconds per $\lambda$; the baseline keeps the 30 seconds of the original runs, which
does not bind on the convex branch.

\subsection{The opponent model}
\label{sec:s2-opponent}

\paragraph{The league baseline and the manager fits.} The league Dirichlet-multinomial model
(3.2)--(3.3) was refitted for every gameweek from 6 to 38 and every position, 132 fits. All
converged: maximum $\hat R=1.01$, minimum bulk effective sample size 1{,}631, and no divergent
transitions. The manager-level update (3.4)--(3.5) was fitted for each gameweek's opponent, on
between 2 and 34 weeks of that manager's history, again 132 fits (maximum $\hat R=1.01$, minimum
bulk effective sample size 1{,}718). Of these, 28 had divergent transitions, 63 in all and at
most 6 in one fit; 17 of the 28 are defender fits. With effective sample sizes this high, a
handful of divergences does not invalidate the posterior means we use, but it marks the
defender fits as the least well-behaved.
The retention coefficient $\gamma_m$ is large throughout: its posterior mean ranges from 2.32 to
10.07, with median 5.3, so a player the manager already owns has his Dirichlet weight
multiplied by $e^{\gamma_m}$, a factor of at least 10. The model has learned what transfer rules
and habit imply, that managers mostly keep their squads.

\paragraph{Overhaul weeks.} A week in which a manager replaces 8 or more of the 15 players is
treated as a Wildcard or Free Hit and is excluded from that manager's history. League-wide, 77
manager-weeks were detected, involving all 20 managers. Most come in clusters that are
implausible as individual decisions: in gameweeks 35 and 36, 19 of the 20 managers are flagged.
The league data show why. Every manager's recorded squad of gameweek 34 shares a median of 2 of
15 players with his gameweek 33 squad and 3 with his gameweek 35 squad, while the gameweek 33 and
35 squads share a median of 13.5. The gameweek 34 squads are therefore a recording error in
the league data, not chips. Similar league-wide breaks occur between gameweeks 5 and 6 and
between 29 and 31, where the median overlap of consecutive squads falls to 4.5 and 3 of 15; we
cannot tell from the data whether these are errors or real. The overhaul filter removes all
of these weeks from the histories, so the selection model is not trained on them.

\subsection{The sparse C-vine}
\label{sec:s2-vine}

Section~3.2.2 specifies a C-vine over the pool of $M$ players (here the 30 most likely
players per position, $M=120$), truncated at tree 15, with each pair-copula chosen by AIC from
six families. The implementation uses a sparse vine instead: Gaussian pair-copulas only, chosen
by BIC, a single tree, pairs with $|\tau|$ below 0.35 set to independence, and only the players
who appear in the manager's history (between 11 and 32 per week) entering the vine. To show why,
we fitted both vines in every gameweek from 6 to 38, on the opponent's latent uniforms $U$
(eqs.~(3.6)--(3.7)) and on a copy of $U$ with each column permuted independently, which destroys
any dependence between players while keeping the marginals. We also recorded the mean held-out
log-likelihood of the last 4, 6 and 8 weeks under a vine fitted on the earlier weeks;
independence scores exactly 0 on this measure. Table~\ref{tab:s2-vine} summarises the results.

\begin{table}[ht]
  \centering
  \caption{Stage 2: vine diagnostics over the 33 gameweeks of the rolling season. Held-out
  log-likelihoods are available in 27 gameweeks (at least five weeks of history must remain).}
  \label{tab:s2-vine}
  \begin{tabular}{lrr}
    \toprule
    & Chapter 3 vine & Sparse vine \\
    \midrule
    Dimension                                   & 120 & 11--32 \\
    Non-independent pairs, real $U$ (mean)      & 378.0 & 0.09 \\
    Non-independent pairs, permuted $U$ (mean)  & 383.2 & --- \\
    Gameweeks with more pairs on permuted $U$   & 22 & --- \\
    Log-likelihood, real $U$ (mean)             & 762.7 & 0.20 \\
    Log-likelihood, permuted $U$ (mean)         & 766.9 & 0.21 \\
    Held-out log-likelihood $<0$ (gameweeks)    & 22 of 27 & 4 of 27 \\
    Held-out log-likelihood, median             & $-957.9$ & 0.0 \\
    \bottomrule
  \end{tabular}
\end{table}

\paragraph{The Chapter 3 vine fits noise.} With the six families and AIC over 120 players, the
vine finds as much dependence in the permuted data as in the real data: on average 378 of its
1{,}680 pairs are non-independent on the real $U$ and 383 on the permuted $U$, with more on the
permuted data in 22 of the 33 gameweeks, and its in-sample log-likelihood is no higher on the
real data (762.7 against 766.9). On held-out weeks it scores below independence in 22 of the 27
gameweeks where the measure is defined, with a median of $-957.9$, and the held-out
log-likelihood is $-\infty$ in 7 of them: the fitted pair-copulas put zero density where later
weeks fall. With a manager history of 4 to 33 weeks and 120 players, AIC has too little data to
reject the flexible families, and the fitted dependence is overfitting. In the first seven
gameweeks, with 4 to 9 weeks of history, every pair falls back to independence.

\paragraph{The sparse vine reduces to independence.} The sparse vine fits no non-independent
pair in 30 of the 33 gameweeks and one pair in each of the other three. Its log-likelihood is
not distinguishable from its permutation null (mean 0.20 against 0.21), it exceeds the null in
only 2 gameweeks, and its held-out log-likelihood is zero in 22 gameweeks, negative in 4 and
positive in 1 (0.17). The data therefore contain no dependence between an opponent's player
choices that a vine can detect beyond the marginals. This is a finding about the data, not a
failure of the software: a manager's history is short, most players are held unchanged for
weeks (so their latent uniforms are draws within one interval of (3.7)), and the transfers that
do occur are too few to estimate pairwise dependence among dozens of players. For §3.2.2 this
means that, with one season of one manager's squads, the copula step adds nothing to the
marginals of §3.2.1. The squads in $\mathbb{W}_o$ are, in effect, drawn from independent
positional marginals and then filtered by the squad rules, the opponent's free transfers and
squad value. The dependence that real squads do have, through the budget, the club limit and
the transfer limit, enters through that filter rather than through the copula.

\paragraph{The size of $\mathbb{W}_o$.} Of the 20{,}000 draws per gameweek, between 7{,}052 and
18{,}327 were feasible squads, and these contained between 124 and 10{,}406 distinct squads
(median 2{,}751). $\mathbb{W}_o$ is small in weeks when the opponent had a single free transfer,
since most feasible draws then repeat the current squad with one change.

\subsection{A worked gameweek}
\label{sec:s2-worked}

To show how one decision is made, we take the gameweek of the $r=50\%$ season in which a
positive $\lambda^*$ gained most in estimated $\mathrm{Prob}\{Y_w>0\}$ over $\lambda=0$. This is
gameweek 9, against manager 593532. Of the 20{,}000 opponent draws, 15{,}959 were feasible
squads, and these contained 1{,}763 distinct squads. Over the Monte Carlo draws the benchmark
$G^{(r')}$ had mean 61.76 and standard deviation 6.68 points. Our current squad's predicted
mean exceeded $\mu_{G^{(r')}}$ by 4.86 points, so the week was on the convex branch of (3.17).

\begin{figure}[ht]
  \centering
  \includegraphics[width=\textwidth]{figures/s2_frontier.pdf}
  \caption{Stage 2, gameweek 9 at $r=50\%$: (a) estimated $\mathrm{Prob}\{Y_{w(\lambda)}>0\}$
  over $\Lambda$, with $\lambda^*$ marked; (b) predicted mean against standard deviation of the
  teams $w(\lambda)$.}
  \label{fig:s2-frontier}
\end{figure}

Figure~\ref{fig:s2-frontier} shows the frontier. The grid produced only four distinct teams. At
$\lambda=0$ the team had mean 68.83 and standard deviation 6.85 points, with estimated
$\mathrm{Prob}\{Y_w>0\}=0.798$. Every $\lambda$ from 0.01 to 0.63 gave a second team, with mean
68.69, standard deviation 8.58 and estimated probability 0.835. Larger values gave weaker
teams: mean 61.12 and probability 0.466 at $\lambda=1.26$ and $2.51$, and mean 56.98 and
probability 0.207 at $\lambda=5$. MVT therefore chose $\lambda^*=0.01$ and made one transfer.

The chosen team has a higher variance of its own than the $\lambda=0$ team, although the week
is on the convex branch. This is consistent with (3.17), because on the convex branch $\lambda$
penalises the variance of the difference $Y_w$, not of our team. That variance falls when our
points covary with $G^{(r')}$, that is, when we hold players the opponent is likely to hold.
For the $\lambda=0.01$ team to be preferred, its covariance with $G^{(r')}$ must have risen by
more than its own variance. In the H2H setting, MVT ahead of the benchmark protects its lead
by matching the opponent rather than by choosing low-variance players. The protection did not
pay off in this week: MVT scored 66 points and the baseline 70, and the opponent's 73 beat both.

\subsection{MVT against the baseline over the season}
\label{sec:s2-season}

We first compare the two optimizers at $r=50\%$. There the benchmark is the median of the
opponent's simulated squads, so $\mathrm{Prob}\{Y_w>0\}$ is the closest analogue of the
probability of winning the match. Both optimizers solved all 33 gameweeks. The $\lambda=0$
problem was solved to optimality in every week. Of the 363 MIQP solves of the MVT season, 166
stopped at the 60-second limit, all of them at $\lambda>0$. Figure~\ref{fig:s2-season} shows
the cumulative points of both optimizers, our real team and its real opponents.

\begin{figure}[ht]
  \centering
  \includegraphics[width=\textwidth]{figures/s2_season.pdf}
  \caption{Stage 2: cumulative points from gameweek 6 to 38 of MVT ($r=50\%$), the baseline,
  our real team and its real opponents, all scored with the same line-up rule.}
  \label{fig:s2-season}
\end{figure}

\paragraph{Points, wins and xP.} MVT scored 1{,}833 points and won 12.5 matches, counting ties
as half a win. The baseline scored 1{,}858 points and won 11. The paired difference in points
is $-0.76$ per gameweek, with a 95\% bootstrap interval of $[-4.36, 3.24]$ (Wilcoxon $p=0.35$,
$n=27$ gameweeks with a non-zero difference). The two optimizers' match results differ in only
4 gameweeks: MVT won 3 of them and the baseline 1 (sign test $p=0.63$). MVT's teams had lower
predicted xP: 67.69 points per gameweek against 71.80, over the 30 gameweeks with an xP target.
Both optimizers made one transfer in every gameweek, and MVT never held. Neither optimizer
comes near our real team. Scored with the same line-up rule, it made 2{,}027 points and won 21
matches, while its real opponents made 1{,}950.

\paragraph{The choice of $\lambda^*$.} In 23 of the 33 gameweeks our current squad's predicted
mean was below $\mu_{G^{(r')}}$, so the week was on the nonconvex branch even at $r=50\%$. The
real opponents' squads were, on the predictions, mostly stronger than ours. $\lambda^*$ was 0
in 23 gameweeks, 0.01 in 7, 0.079 in 2 and 0.158 in 1. Of the ten positive values, six were on
the convex branch and four on the nonconvex branch. The chosen team's standard deviation
averaged 10.53 points, against the baseline's 9.77.

\paragraph{Calibration.} MVT's mean estimated $\mathrm{Prob}\{Y_w>0\}$ was 0.62, against an
observed win rate of 0.38 (Brier score 0.382). The estimates are therefore too high, as in
Stage~1, though by less: the opponent model describes a single real manager, while the
simulated field of Stage~1 was weaker than the real contest. The upward bias of $\mu_\delta$ for
the top-predicted players (Section~\ref{sec:predictions}) applies here as well.

\subsection{The effect of $r$}
\label{sec:s2-r}

The MVT season was run at $r\in\{50,60,80,90,95,99,100\}\%$ against the same baseline.
Table~\ref{tab:s2-r} and Figure~\ref{fig:s2-r} report the results. Every season solved all 33
gameweeks, and no gameweek needed the single-squad fallback. The $\lambda=0$ problem was solved
to optimality in every gameweek except one: at $r=99\%$, gameweek 38 stopped at the time limit
and kept its best team found. Between 145 and 187 of each season's 363 solves stopped at the
60-second limit, so at $\lambda>0$ the frontier consists of the best teams found, not proven
optima. At $r=100\%$ the benchmark equalled the maximum over $\mathbb{W}_o$ in every draw of
every gameweek, as intended.

\begin{table}[ht]
  \centering
  \caption{Stage 2: MVT at each $r$ against the baseline, 33 gameweeks. Nonconvex: gameweeks on
  the nonconvex branch of (3.17). $\sigma_w$: mean standard deviation of the chosen team's points
  under $\Sigma_\delta$. $\Delta$ points: mean difference per gameweek (MVT minus baseline) with a
  95\% bootstrap interval; $p$: Wilcoxon signed-rank test. Wins count ties as 0.5; $p_W$: sign
  test on the gameweeks where the two match results differ. $\hat P$: mean estimated
  $\mathrm{Prob}\{Y_w>0\}$.}
  \label{tab:s2-r}
  \footnotesize
  \setlength{\tabcolsep}{3pt}
  \begin{tabular}{rrrrrrrrrrr}
    \toprule
    $r$ & Nonconvex & $\lambda^*\neq0$ & Held & $\sigma_w$ & Points & $\Delta$ points [95\% CI] & $p$ & Wins & $p_W$ & $\hat P$ \\
    \midrule
    50\%  & 23 & 10 & 0 & 10.53 & 1833 & $-0.76$ [$-4.36$, $3.24$] & 0.348 & 12.5 & 0.625 & 0.622 \\
    60\%  & 21 & 13 & 0 &  9.30 & 1822 & $-1.09$ [$-5.42$, $3.12$] & 0.559 & 12.5 & 0.625 & 0.550 \\
    80\%  & 22 &  9 & 0 & 10.27 & 1881 & $+0.70$ [$-2.09$, $3.67$] & 0.880 & 13.5 & 0.250 & 0.557 \\
    90\%  & 27 & 11 & 0 & 10.23 & 1891 & $+1.00$ [$-1.97$, $4.21$] & 0.776 & 13.0 & 0.250 & 0.478 \\
    95\%  & 28 & 13 & 0 & 10.33 & 1902 & $+1.33$ [$-1.48$, $4.30$] & 0.546 & 12.5 & 0.500 & 0.430 \\
    99\%  & 33 & 20 & 0 & 12.12 & 1775 & $-2.52$ [$-7.48$, $2.52$] & 0.335 & 16.0 & 0.039 & 0.247 \\
    100\% & 33 & 18 & 5 & 11.56 & 1702 & $-4.73$ [$-8.45$, $-1.00$] & 0.017 & 10.0 & 1.000 & 0.158 \\
    \midrule
    Baseline & 0 & --- & 0 & 9.77 & 1858 & --- & --- & 11.0 & --- & --- \\
    \bottomrule
  \end{tabular}
\end{table}

\begin{figure}[ht]
  \centering
  \includegraphics[width=\textwidth]{figures/s2_r_effect.pdf}
  \caption{Stage 2: the effect of $r$. (a) Season points; (b) H2H win rate; (c) mean standard
  deviation of the chosen team's points; (d) MVT's mean estimated $\mathrm{Prob}\{Y_w>0\}$ and
  observed win rate.}
  \label{fig:s2-r}
\end{figure}

\paragraph{The branch and the variance of the team.} As in Stage~1, the branch of (3.17)
follows $r$. Between 21 and 23 gameweeks were on the nonconvex branch for $r\le 80\%$, 27 and
28 at $r=90\%$ and $95\%$, and all 33 at $r=99\%$ and $100\%$. Unlike Stage~1, most weeks are on
the nonconvex branch even at low $r$, because the opponents' squads were predicted to be
stronger than ours. $\lambda^*$ was 0 in 20 to 24 gameweeks for $r\le 95\%$ and positive in 20
and 18 at $r=99\%$ and $100\%$. At $r=100\%$ MVT also held its squad in 5 gameweeks, where no
team improved on holding by the 0.03 margin. The chosen team's standard deviation is close to
the baseline's for $r\le 95\%$ (9.30 to 10.53 against 9.77) and higher at $r=99\%$ and $100\%$
(12.12 and 11.56). MVT does not reduce the team's variance at low $r$ as it did in Stage~1, for
two reasons. Most weeks are on the nonconvex branch. On the convex branch, the variance that
(3.17) penalises is that of the difference with the opponent, which can fall while our team's
own variance rises (Section~\ref{sec:s2-worked}).

\paragraph{Points.} No value of $r$ gives MVT significantly more points than the baseline. The
differences for $r\le 95\%$ lie between $-1.09$ and $+1.33$ points per gameweek, with
Wilcoxon $p\ge 0.35$. At $r=100\%$ MVT scored 4.73 points per gameweek fewer ($p=0.017$, $n=31$),
and at $r=99\%$ 2.52 fewer ($p=0.34$). As in Stage~1, the variance bought on the nonconvex
branch costs mean points. Like Stage~1's significant results, $p=0.017$ does not survive a
Bonferroni correction for the seven values of $r$ (threshold 0.007).

\paragraph{Wins.} MVT won between 12.5 and 13.5 matches for $r\le 95\%$, against the baseline's
11. Its match results differed from the baseline's in at most 3 gameweeks, so none of these
differences can be significant. The most wins came at $r=99\%$: 16 against 11, with MVT winning
8 of the 9 discordant gameweeks (sign test $p=0.039$), despite scoring 83 fewer points over the
season. At $r=100\%$ MVT won 10. A higher-variance team can win more often while scoring fewer
points in total, and (3.17) aims at exactly this trade-off. However, the $r=99\%$ result does not
survive the Bonferroni correction, the neighbouring values of $r$ do not show it, and one season
of 33 matches cannot separate it from chance.

\paragraph{Calibration.} For $r>50\%$ the benchmark lies above the median of the opponent's
squads, so $\mathrm{Prob}\{Y_w>0\}$ is a probability of beating a stronger squad than the
opponent's typical one, not of winning the match. The mean estimate falls with $r$, from 0.62
to 0.16, while the observed win rate stays between 0.30 and 0.48. The estimate is above the win
rate for $r\le 95\%$ and below it at $r=99\%$ and $100\%$, as the stricter benchmark implies.

\section{Discussion}
\label{sec:discussion}

\paragraph{MVT against the baseline.} The two stages agree on how MVT behaves and on what it
achieves. In both, the branch of (3.17) follows $r$, and MVT buys variance when it expects to be
behind the benchmark. In both, that variance costs mean points at $r=100\%$: $-5.24$ points per
gameweek in Stage~1 and $-4.73$ in Stage~2. In neither stage does MVT score significantly more
points than the baseline at any $r$, or hit its target significantly more often once the
seven values of $r$ are accounted for. MVT's estimated probabilities are too high at $r=50\%$ in
both stages, far more so in Stage~1 (0.996 against an observed 0.52) than in Stage~2 (0.62
against 0.38). With 33 gameweeks per stage, the method's advantage over a points maximiser, if
it exists, is smaller than one season can detect.

\paragraph{Winner-takes-all against lower $r$.} In a FantasyGo contest, $r$ is set by the
payout: a team must finish in the top $1-r$ of a field of hundreds to be paid, and the field's
quantile is the natural benchmark. At low $r$ our team is usually ahead of that benchmark, and
MVT buys safety; in a winner-takes-all contest ($r=100\%$) it must beat the best entry, is
always behind, and buys variance. An H2H match has no field. The match is won by beating one
opponent, and the benchmark that represents this most directly is $r=50\%$, the median of the
opponent's possible squads. A higher $r$ asks MVT to beat the opponent's best plausible squad,
a stricter target than the match requires. It moves every week to the nonconvex branch, and at
$r=100\%$ it cost points without adding wins. In H2H, $r$ is therefore not dictated by the
rules but acts as a risk setting. The single season suggests that a high setting ($r=99\%$) may
convert variance into wins, but the evidence is not strong enough to recommend it.

\paragraph{Non-strategic opponents.} The opponents in both stages are modelled from their own
past behaviour, and our choice does not enter $\mathbb{W}_o$. In a FantasyGo contest with
hundreds of entries this is close to true: no single entry can track and respond to ours. In an
H2H league it is an approximation. Each manager knows his opponent in advance and can see that
opponent's squad from previous gameweeks. A strategic opponent could respond to our expected
squad, for example by doubling our likely captain to cancel his points or by avoiding our
players to create variance. The model would then misstate both the location of $G^{(r')}$ and
its covariance $\sigma_{\delta,G^{(r')}}$ with our players, which is the term through which
(3.17) decides whether to copy or avoid the opponent's players. The stability of the squads we
observe ($\gamma_m$ between 2.32 and 10.07) suggests that, in this league, managers mostly
manage their own squads rather than respond to opponents. A game-theoretic treatment, as in the
opponent classification and the ``sword'' and ``shield'' strategies outlined in §1.4, is
outside what we implemented.

\paragraph{The objectives of §1.3.2.} (i) The H2H problem was formulated as the mean-variance
problem of Chapter~3, with the opponent's simulated squads as the benchmark, and solved as a
MIQP over $\Lambda$ in every gameweek of every season; the $\lambda=0$ problem was solved to
optimality in all but one gameweek, and the problems at $\lambda>0$ often to the time limit.
(ii) Player points were predicted per position (Section~\ref{sec:predictions}), with the
caveat that the predictions overstate the points of the top-predicted players. (iii) The
opponents' squads were modelled by the league and manager selection models, which converged;
the copula step reduced to independence, so the squads are drawn from the positional
marginals filtered by the squad rules. (iv) The rolling season produced a transfer, starting
XI and captain in every gameweek. (v) On 2024-25 FPL data the method showed no significant
advantage over the baseline in points or wins. Both optimizers stayed well below our real
team (2{,}027 points and 21 wins, against between 1{,}702 and 1{,}902 points and between 10 and
16 wins), which points to the predictions, rather than the optimization, as the main
limitation.

\paragraph{Limitations.} Each stage is a single season of 33 gameweeks, so paired tests have
little power, and a handful of contests decide the prize totals. $\Sigma_\delta$ rests on at
most 33 gameweeks of errors for several hundred players; its correlations are noisy, and the
eigenvalue floor keeps it usable without making it accurate. $\mu_\delta$ overstates the
realised points of the top-predicted players by about one point per player
(Section~\ref{sec:predictions}), which makes every estimated probability optimistic. The
opponent model does not include chips: Wildcard and Free Hit weeks are excluded from the
histories, but our model does not predict when an opponent will play one. The league data
contain an erroneous gameweek of squads (gameweek 34) and two league-wide breaks of unknown
origin. In Stage~2 the two optimizers' match results differ in at most 9 of the 33 gameweeks,
so every comparison of wins rests on a handful of matches. Between 145 and 187 of each MVT
season's 363 solves stopped at the 60-second limit, nearly all at $\lambda>0$, so those teams
are the best found rather than proven optima. Because the limit is in wall-clock time, a run on a loaded machine can return
different teams; the baseline reported here was run alone and reproduces the original
baseline's decisions exactly. Finally, our line-ups and those of the opponents and our real
team are all chosen by the same rule, so the comparison is fair, but the real managers'
line-up decisions are not modelled.
